\documentclass[12pt,letterpaper]{article}
\usepackage{graphicx,textcomp}
\usepackage{natbib}
\usepackage{setspace}
\usepackage{fullpage}
\usepackage{color}
\usepackage[reqno]{amsmath}
\usepackage{amsthm}
\usepackage{fancyvrb}
\usepackage{amssymb,enumerate}
\usepackage[all]{xy}
\usepackage{endnotes}
\usepackage{lscape}
\newtheorem{com}{Comment}
\usepackage{float}
\usepackage{hyperref}
\newtheorem{lem} {Lemma}
\newtheorem{prop}{Proposition}
\newtheorem{thm}{Theorem}
\newtheorem{defn}{Definition}
\newtheorem{cor}{Corollary}
\newtheorem{obs}{Observation}
\usepackage[compact]{titlesec}
\usepackage{dcolumn}
\usepackage{tikz}
\usetikzlibrary{arrows}
\usepackage{multirow}
\usepackage{xcolor}
\newcolumntype{.}{D{.}{.}{-1}}
\newcolumntype{d}[1]{D{.}{.}{#1}}
\definecolor{light-gray}{gray}{0.65}
\usepackage{url}
\usepackage{listings}
\usepackage{color}
\usepackage{caption}

\definecolor{codegreen}{rgb}{0,0.6,0}
\definecolor{codegray}{rgb}{0.5,0.5,0.5}
\definecolor{codepurple}{rgb}{0.58,0,0.82}
\definecolor{backcolour}{rgb}{0.95,0.95,0.92}

\lstdefinestyle{mystyle}{
	backgroundcolor=\color{backcolour},   
	commentstyle=\color{codegreen},
	keywordstyle=\color{magenta},
	numberstyle=\tiny\color{codegray},
	stringstyle=\color{codepurple},
	basicstyle=\footnotesize,
	breakatwhitespace=false,         
	breaklines=true,                 
	captionpos=b,                    
	keepspaces=true,                 
	numbers=left,                    
	numbersep=5pt,                  
	showspaces=false,                
	showstringspaces=false,
	showtabs=false,                  
	tabsize=2
}
\lstset{style=mystyle}
\newcommand{\Sref}[1]{Section~\ref{#1}}
\newtheorem{hyp}{Hypothesis}


\title{Problem Set 4}
\date{Due: November 27, 2025}
\author{Applied Stats/Quant Methods 1\\
	Shelly Veal-Upham\\
	25337422}

\begin{document}
	\maketitle
	\section*{Instructions}
	\begin{itemize}
		\item This problem set is due before 23:59 on Thursday November 27, 2025. No late assignments will be accepted.
	\end{itemize}

	\vspace{.5cm}
\section*{Question 1: Economics}
\vspace{.25cm}
\noindent 	
In this question, use the \texttt{prestige} dataset in the \texttt{car} library. First, run the following commands:

\begin{verbatim}
install.packages(car)
library(car)
data(Prestige)
help(Prestige)
\end{verbatim} 


\noindent We would like to study whether individuals with higher levels of income have more prestigious jobs. Moreover, we would like to study whether professionals have more prestigious jobs than blue and white collar workers.\\

\vspace{.1cm}
*Note that I've renamed the dataframe "df".

\newpage
\begin{enumerate}
	
	\item [(a)]
	Create a new variable \texttt{professional} by recoding the variable \texttt{type} so that professionals are coded as $1$, and blue and white collar workers are coded as $0$ (Hint: \texttt{ifelse}).\\
	
	\vspace{.1cm}
		\lstinputlisting[language=R, firstline=44, lastline=44]{PS04_SVU.R} 
	\vspace{.1cm}

	\item [(b)]
	Run a linear model with \texttt{prestige} as an outcome and \texttt{income}, \texttt{professional}, and the interaction of the two as predictors (Note: this is a continuous $\times$ dummy interaction.)\\
	\vspace{.1cm}
		\lstinputlisting[language=R, firstline=79, lastline=79]{PS04_SVU.R} 
	\vspace{.2cm}
\begin{table}[!htbp] \centering 
	\caption{Impact of Income and Professionality on Prestige} 
	\label{} 
	\begin{tabular}{@{\extracolsep{5pt}}lc} 
		\\[-1.8ex]\hline 
		\hline \\[-1.8ex] 
		& \multicolumn{1}{c}{\textit{Dependent variable:}} \\ 
		\cline{2-2} 
		\\[-1.8ex] & Prestige \\ 
		\hline \\[-1.8ex] 
		Income & 0.003$^{***}$ \\ 
		& (0.0005) \\ 
		& \\ 
		Professional & 37.781$^{***}$ \\ 
		& (4.248) \\ 
		& \\ 
		Interaction & $-$0.002$^{***}$ \\ 
		& (0.001) \\ 
		& \\ 
		Constant & 21.142$^{***}$ \\ 
		& (2.804) \\ 
		& \\ 
		\hline \\[-1.8ex] 
		Observations & 98 \\ 
		R$^{2}$ & 0.787 \\ 
		Adjusted R$^{2}$ & 0.780 \\ 
		Residual Std. Error & 8.012 (df = 94) \\ 
		F Statistic & 115.878$^{***}$ (df = 3; 94) \\ 
		\hline 
		\hline \\[-1.8ex] 
		\textit{Note:}  & \multicolumn{1}{r}{$^{*}$p$<$0.1; $^{**}$p$<$0.05; $^{***}$p$<$0.01} \\ 
	\end{tabular} 
\end{table}
\noindent For visual verification of these impacts see plot below, noting that the slope of the red line is equal to 0.00317 ($\beta{1}$), and the green line slope is equal to 0.000845 ($\beta_{1} - \beta_{3}$):\\
	\begin{figure}[h!]\centering
		\caption{\footnotesize Income vs Prestige.}
		\label{fig:plot_1}
		\includegraphics[width=.85\textwidth]{plot_1.pdf}
	\end{figure}
\newpage

The code for creating this plot can be seen here:
		\lstinputlisting[language=R, firstline=49, lastline=62]{PS04_SVU.R} 
	\item [(c)]
	Write the prediction equation based on the result.\\
	
	\vspace{.08cm}
		\(\widehat{\texttt{prestige}} = 21.142 + 0.00317\texttt{income} + 37.781\texttt{professional} - 0.00233\texttt{income*professional}\)
	\vspace{.1cm}
	
	\item [(d)]
	Interpret the coefficient for \texttt{income}.\\
	
	\vspace{.1cm}
		Income's coefficient (0.00317) is rather small, but this makes sense given that income is measured in dollars. Interpreting it as is, a one-dollar gain in income results in an increase of 0.00317 points on the prestige scale for non-professional jobs (professional = 0). It is worth noting, however, that a significant change in income would be in the hundreds if not thousands of dollars, as we will explore in item f. \\
	\vspace{.1cm}

	\item [(e)]
	Interpret the coefficient for \texttt{professional}.\\
	
	\vspace{.1cm}
		Professional's coefficient, 37.781, is much larger than income's coefficient, and it being a dummy variable for professionality (as opposed to white or blue collar) means we interpret it as follows: professional jobs enjoy 37.781 points more, on average, in prestige compared to white or blue collar jobs when holding income constant. \\
	\vspace{.1cm}
	
	\item [(f)]
	What is the effect of a \$1,000 increase in income on prestige score for professional occupations? In other words, we are interested in the marginal effect of income when the variable \texttt{professional} takes the value of $1$. Calculate the change in $\hat{y}$ associated with a \$1,000 increase in income based on your answer for (c).\\

	\vspace{.1cm}
		A \$1,000 increase in income for professional occupations will, on average, result in a 0.845 point gain on the prestige scale. For example, professionals with an income of \$10,000 enjoy a prestige score of approximately 67.323. A \$1,000 increase in income would move them to a prestige score of approximately 68.163. Note in the code below: b1 + b3 results in the change in income for professional occupations (0.000845), which for \$1,000 increase in income results in a 0.845 point gain in prestige. \\
				\lstinputlisting[language=R, firstline=111, lastline=120]{PS04_SVU.R} 
				\lstinputlisting[language=R, firstline=127, lastline=130]{PS04_SVU.R} 
	\vspace{.1cm}
	
	\item [(g)]
	What is the effect of changing one's occupations from non-professional to professional when her income is \$6,000? We are interested in the marginal effect of professional jobs when the variable \texttt{income} takes the value of $6,000$. Calculate the change in $\hat{y}$ based on your answer for (c).\\
	
	\vspace{.1cm}
		It's first noticeable here that the question phrasing is a bit odd given the data -- one wouldn't magically change their profession and hold their income constant. Holding income at \$6,000 really shows us how people in different profession types receive different prestige scores even if their income is the same. As such, professionals earning \$6,000 on average enjoy 23.827 more points on the prestige scale than their white collar or blue collar counterparts who earn the same income (63.995 for professionals, 40.168 for non-professionals).\\
			In R (\texttt{ans1g} = 23.827):
					\lstinputlisting[language=R, firstline=133, lastline=134]{PS04_SVU.R} 
	\vspace{.1cm}
	
	
\end{enumerate}

\newpage

\section*{Question 2: Political Science}
\vspace{.25cm}
\noindent 	Researchers are interested in learning the effect of all of those yard signs on voting preferences.\footnote{Donald P. Green, Jonathan	S. Krasno, Alexander Coppock, Benjamin D. Farrer,	Brandon Lenoir, Joshua N. Zingher. 2016. ``The effects of lawn signs on vote outcomes: Results from four randomized field experiments.'' Electoral Studies 41: 143-150. } Working with a campaign in Fairfax County, Virginia, 131 precincts were randomly divided into a treatment and control group. In 30 precincts, signs were posted around the precinct that read, ``For Sale: Terry McAuliffe. Don't Sellout Virgina on November 5.'' \\

Below is the result of a regression with two variables and a constant.  The dependent variable is the proportion of the vote that went to McAuliff's opponent Ken Cuccinelli. The first variable indicates whether a precinct was randomly assigned to have the sign against McAuliffe posted. The second variable indicates
a precinct that was adjacent to a precinct in the treatment group (since people in those precincts might be exposed to the signs).  \\

\vspace{.5cm}
\begin{table}[!htbp]
	\centering 
	\textbf{Impact of lawn signs on vote share}\\
	\begin{tabular}{@{\extracolsep{5pt}}lccc} 
		\\[-1.8ex] 
		\hline \\[-1.8ex]
		Precinct assigned lawn signs  (n=30)  & 0.042\\
		& (0.016) \\
		Precinct adjacent to lawn signs (n=76) & 0.042 \\
		&  (0.013) \\
		Constant  & 0.302\\
		& (0.011)
		\\
		\hline \\
	\end{tabular}\\
	\footnotesize{\textit{Notes:} $R^2$=0.094, N=131}
\end{table}

\vspace{.5cm}
\begin{enumerate}
	\item [(a)] Use the results from a linear regression to determine whether having these yard signs in a precinct affects vote share (e.g., conduct a hypothesis test with $\alpha = .05$).\\
	\vspace{.05cm}
		\begin{center}
			\(t_{assigned} = \frac{0.042 - 0}{0.016} = 2.625\) \\
			Yields p-value of 0.00972\\
		\vspace{.05cm}
		\end{center}
	\newpage
	In R:\\
			\lstinputlisting[language=R, firstline=145, lastline=161]{PS04_SVU.R} 
	\vspace{.03cm}		
	With $\alpha = .05$, a p-value of 0.00972 indicates significance, suggesting that vote share \textbf{is} impacted by the postage of signs in voting precincts. In other words, there is a significant difference between the reference category (no signage, not adjacent to signed precincts) and this category (signage posted). \\
	
	\item [(b)]  Use the results to determine whether being
	next to precincts with these yard signs affects vote
	share (e.g., conduct a hypothesis test with $\alpha = .05$).
	\vspace{.05cm}
		\begin{center}
			\(t_{adjacent} = \frac{0.042 - 0}{0.013} = 3.231\) \\
			Yields p-value of 0.00157\\
		\vspace{.05cm}
		\end{center}
	In R:\\
	\lstinputlisting[language=R, firstline=165, lastline=167]{PS04_SVU.R} 
	\vspace{.03cm}		
	With $\alpha = .05$, a p-value of 0.00157 indicates significance, suggesting that vote share \textbf{is} affected by adjacence to precincts with signs posted.\\
	\item [(c)] Interpret the coefficient for the constant term substantively.\\
	
	\vspace{.05cm}
		The constant term coefficient, 0.302, indicates our estimated proportion of the vote for precincts without interaction with the signs posted (not having signs posted, and not being adjacent to precincts with posted signs -- our reference category). This suggests that Cuccinelli is receiving, on average, 30 percent of the vote share in precincts without the influence of the signage.\\
	\vspace{.05cm}
	\item [(d)] Evaluate the model fit for this regression.  What does this	tell us about the importance of yard signs versus other factors that are not modeled?
	
	\vspace{.05cm}
		The $R^2$ value of 0.094 suggests that the model explains only approximately 9 percent of the variation in the data. Indeed, we would assume that there are other factors that influence voting much more than posted signage, such as incumbency status, the state of the economy, demographic factors of the voting population and the candidate, etc.\\
	\vspace{.05cm}
\end{enumerate}  


\end{document}
